% Generated draft: requires mathematical and source review.
\documentclass[11pt]{article}
\usepackage[margin=25mm]{geometry}
\usepackage{fontspec}
\setmainfont{Latin Modern Roman}
\usepackage{amsmath,amssymb,mathtools,unicode-math}
\setmathfont{Latin Modern Math}
\usepackage{enumitem,needspace,xurl,hyperref}
\setlength{\parindent}{0pt}
\setlength{\parskip}{5pt}
\setlength{\emergencystretch}{3em}
\newcommand{\N}{\mathbb{N}}
\newcommand{\Z}{\mathbb{Z}}
\newcommand{\Q}{\mathbb{Q}}
\newcommand{\R}{\mathbb{R}}
\newcommand{\C}{\mathbb{C}}
\newcommand{\F}{\mathbb{F}}
\newcommand{\A}{\mathbb{A}}
\newcommand{\PP}{\mathbb{P}}
\newcommand{\E}{\mathbb{E}}
\newcommand{\eps}{\varepsilon}
\newcommand{\dd}{\,\mathrm{d}}
\begin{document}
\section*{UnsolvedMath ID: 10600032}
\textbf{Source problem code:} catalogue entry
\subsection{Definitions and mathematical statement}
An ordinary two-knot in four-space is an embedding \(f:S^2\hookrightarrow\R^4\), considered up to ambient isotopy in a specified category. The source asks to generalize virtual knot theory to \emph{virtual two-spheres} in four-space. A precise proposed theory must specify its generalized objects and their equivalence, how ordinary embedded two-spheres are represented, and how it extends the virtual-crossing/stabilization idea from one-dimensional knots.

For a diagrammatic approach, one would need to identify the admissible projected surface data and any virtual local singularity labels, give the complete allowed local moves, and state how ordinary surface-diagram moves are incorporated. For a geometric approach, one would need to specify the ambient objects and stabilization/equivalence operations. These are alternative ways to make the construction concrete, not a choice imposed by the source. A resulting comparison with ordinary two-knot isotopy must be stated; the excerpt does not assert injectivity of such a comparison or prescribe a particular higher-dimensional diagram category.

The requested source surface is \(S^2\), not an arbitrary-genus surface. Smooth, locally flat piecewise-linear or topological embedding conventions, orientation and the exact meaning of ``generalize'' are unstated. Consequently this is a programme to define a theory with justified equivalence rules, not an existence theorem about an already specified set of virtual two-spheres.
\subsection{Short English statement}
Develop a virtual analogue of embedded two-spheres in four-space, including precise objects, moves or stabilizations, and their relation to ordinary two-knot isotopy. The source does not prescribe the higher-dimensional virtual model.
\subsection{Sources}
\begin{itemize}
\item Local UnsolvedMath registry, permanent ID \(10600032\), exact source code \texttt{catalogue entry}.
\item \emph{Unsolved Problems in Virtual Knot Theory and Combinatorial Knot Theory}, revised survey, problem-list item 32: \url{https://arxiv.org/abs/1409.2823}. Registry access 2026-07-29; source-tex extraction.
\item The research summary mentions spinning constructions and Gauss-diagram invariants in scattered virtual-surface work, but supplies no specific references.
\end{itemize}
\end{document}
